\documentclass[11pt]{article}
\usepackage[utf8]{inputenc}
\usepackage{amsmath,amsthm,amssymb}
\usepackage{geometry}
\geometry{margin=2.5cm}
\usepackage[hidelinks]{hyperref}
\usepackage{booktabs}

\newtheorem{theorem}{Theorem}
\newtheorem{lemma}[theorem]{Lemma}
\newtheorem{corollary}[theorem]{Corollary}
\newtheorem{proposition}[theorem]{Proposition}
\theoremstyle{definition}
\newtheorem{definition}[theorem]{Definition}
\newtheorem{example}[theorem]{Example}
\theoremstyle{remark}
\newtheorem{remark}[theorem]{Remark}

\DeclareMathOperator{\arccosh}{arccosh}

\title{Rational Square Root Approximation\\via the Pell--Chebyshev Series}
\author{Jan Popelka\thanks{Email: popojan@protonmail.com}}
\date{}

\begin{document}
\maketitle

\begin{abstract}
Given the fundamental solution $(x,y)$ to $x^2 - ny^2 = 1$,
we define a series of rational approximations to~$\sqrt{n}$
whose terms involve Chebyshev polynomials evaluated at~$x$.
We prove that the $k$-th partial sum collapses to a single
hyperbolic expression:
$E_k = \sqrt{n}\,\tanh\bigl(\tfrac{k+1}{2}\,R\bigr)$,
where $R = \log(x + y\sqrt{n})$ is the Pell regulator.
The interval width at depth~$k$ is $\sqrt{n}/\sinh\bigl((k{+}1)\,R\bigr)$,
giving exponential convergence with rate~$1/\varepsilon$ per term,
where $\varepsilon$ is the fundamental unit.
\end{abstract}

%% ====================================================================
\section{The series}
%% ====================================================================

Let $(x,y)$ be the fundamental solution to $x^2 - ny^2 = 1$,
and let $T_m$, $U_m$ denote Chebyshev polynomials of the first
and second kind.

\begin{definition}\label{def:series}
The \emph{Egypt series} for $\sqrt{n}$ at depth~$k$ is
\begin{equation}\label{eq:egypt}
E_k \;=\; \frac{x-1}{y}\,\Bigl(1 + \sum_{j=1}^{k}
  \frac{1}{T_{\lceil j/2\rceil}(x)\bigl(
  U_{\lfloor j/2\rfloor}(x)
  - U_{\lfloor j/2\rfloor - 1}(x)\bigr)}\Bigr).
\end{equation}
\end{definition}

Each summand is the reciprocal of a product of Chebyshev polynomials
at the Pell solution~$x$.
The first few terms are:
\begin{alignat*}{2}
\text{Term}_1 &= \frac{1}{T_1(x) \cdot (U_0(x) - U_{-1}(x))}
              &&= \frac{1}{x}, \\[3pt]
\text{Term}_2 &= \frac{1}{T_1(x) \cdot (U_1(x) - U_0(x))}
              &&= \frac{1}{x(2x-1)}.
\end{alignat*}

\begin{remark}
An equivalent formulation uses factorials (the author's
\texttt{egypt} package~\cite{egypt-repo}):
\[
\text{Term}_j = \frac{1}{1 + \sum_{i=1}^{j}
  \frac{2^{i-1}\,(x{-}1)^i\,(j{+}i)!}{(j{-}i)!\,(2i)!}},
\]
with argument $x - 1$ instead of~$x$.
The equivalence of the Chebyshev and factorial forms has been
verified numerically for $j \le 200$;
an algebraic proof is given in~\cite{Popelka2025egypt}.
\end{remark}

%% ====================================================================
\section{Pell powers and the Egypt identity}
%% ====================================================================

Define the $k$-th Pell power $(x_k, y_k)$ by
$x_k + y_k\sqrt{n} = (x + y\sqrt{n})^k$, i.e.,
$x_k = T_k(x)$, $y_k = y \cdot U_{k-1}(x)$.

\begin{proposition}\label{prop:pellpower}
$E_k = (x_{k+1} - 1)\,/\,y_{k+1}$.
\end{proposition}

\begin{proof}
By induction, using the recurrence for $(x_k, y_k)$ and
the definition of the Egypt terms.
Alternatively, verify that each term telescopes:
the partial sum $1 + \sum_{j=1}^{k} \text{Term}_j$
equals $y \cdot U_k(x) \cdot (T_{k+1}(x) - 1) / ((x-1) \cdot y_{k+1})$.
\end{proof}

This shows that the Egypt series at depth~$k$ produces the
\emph{same} rational bounds as the $(k{+}1)$-th Pell power,
repackaged as a cumulative sum.

%% ====================================================================
\section{The hyperbolic collapse}\label{sec:collapse}
%% ====================================================================

Parametrize the Pell conic by $x = \cosh R$, $y = \sinh R/\sqrt{n}$,
where $R = \arccosh(x) = \log(x + y\sqrt{n})$ is the regulator.

\begin{theorem}[Tanh collapse]\label{thm:tanh}
For all $k \ge 0$:
\begin{equation}\label{eq:tanh}
\boxed{\;E_k \;=\; \sqrt{n}\;\tanh\!\Bigl(\frac{(k+1)\,R}{2}\Bigr).\;}
\end{equation}
\end{theorem}

\begin{proof}
By Proposition~\ref{prop:pellpower}:
\[
E_k = \frac{x_{k+1} - 1}{y_{k+1}}
    = \frac{\cosh\bigl((k{+}1)R\bigr) - 1}
           {\sinh\bigl((k{+}1)R\bigr)/\sqrt{n}}
    = \sqrt{n}\;\frac{\cosh S - 1}{\sinh S},
\]
where $S = (k{+}1)R$.
The identity $(\cosh S - 1)/\sinh S = \tanh(S/2)$
(from $\cosh S - 1 = 2\sinh^2(S/2)$ and
$\sinh S = 2\sinh(S/2)\cosh(S/2)$) completes the proof.
\end{proof}

\begin{corollary}[Bounds and width]\label{cor:bounds}
The upper bound $x_{k+1}/y_{k+1}$ satisfies
\[
\frac{x_{k+1}}{y_{k+1}} = \sqrt{n}\;\coth\bigl((k{+}1)R\bigr)
= \frac{\sqrt{n}}{\tanh\bigl((k{+}1)R\bigr)}.
\]
The interval width is
\begin{equation}\label{eq:width}
\frac{x_{k+1}}{y_{k+1}} - \frac{x_{k+1} - 1}{y_{k+1}}
= \frac{1}{y_{k+1}}
= \frac{\sqrt{n}}{\sinh\bigl((k{+}1)R\bigr)},
\end{equation}
which decreases exponentially:
$\text{width}_k / \text{width}_{k+1} = 2\cosh R = 2x$ for
successive depths.
\end{corollary}

\begin{proof}
The width is $1/y_{k+1} = \sqrt{n}/\sinh((k{+}1)R)$.
The ratio $\sinh((k{+}2)R)/\sinh((k{+}1)R)
= 2\cosh R$ by the identity
$\sinh((k{+}2)R) = 2\cosh R \cdot \sinh((k{+}1)R) - \sinh(kR)$
evaluated at $k \to \infty$; for finite~$k$ the ratio approaches
$2\cosh R = 2x$ monotonically from below.
\end{proof}

\begin{remark}
Each Egypt term improves precision by exactly $R$~nats
($= \log_2\varepsilon$ bits), where $\varepsilon = x + y\sqrt{n}$
is the fundamental unit.
The doubly-exponential convergence reported for the series~\eqref{eq:egypt}
is the same exponential convergence of $\tanh(S/2) \to 1$,
restated in terms of the cumulative sum.
\end{remark}

%% ====================================================================
\section{Symbolic evaluation}
%% ====================================================================

Computer algebra systems can evaluate~\eqref{eq:tanh}
to exact rationals.  Since $\varepsilon = x + y\sqrt{n}$ is algebraic,
\[
E_k = \sqrt{n}\;\frac{\varepsilon^{(k+1)/2} - \varepsilon^{-(k+1)/2}}
     {\varepsilon^{(k+1)/2} + \varepsilon^{-(k+1)/2}},
\]
and the irrationalities cancel by the Pell relation.
In Wolfram Language:

\medskip
\begin{center}
\ttfamily Sqrt[n] * Tanh[(k+1)/2 * Log[x + y*Sqrt[n]]] // FullSimplify
\end{center}
\medskip

\begin{example}
For $n = 13$, $(x,y) = (649, 180)$:
\begin{center}
\begin{tabular}{cll}
\toprule
$k$ & $E_k$ & Decimal \\
\midrule
$0$ & $18/5$     & $3.600\ldots$ \\
$1$ & $2340/649$ & $3.60554\ldots$ \\
$2$ & $23382/6485$ & $3.6055512\ldots$ \\
\bottomrule
\end{tabular}
\end{center}
Compare $\sqrt{13} = 3.60555127\ldots$\;
Each term adds $\approx \log_{10}\varepsilon \approx 2.8$ decimal digits.
\end{example}

%% ====================================================================
\section{Chebyshev composition for solvable families}
%% ====================================================================

When $n$ belongs to a family solvable by the Chebyshev--Pell
method~\cite{PopelkaPellTower2025}
(i.e., $x = T_m(z)$ for a rational seed~$z$ with
$\operatorname{denom}(z) \le 2$), the composition identities
\begin{align}
T_j\bigl(T_m(z)\bigr) &= T_{jm}(z), \label{eq:Tcomp} \\[3pt]
U_{j-1}\bigl(T_m(z)\bigr) &= \frac{U_{jm-1}(z)}{U_{m-1}(z)}
\label{eq:Ucomp}
\end{align}
(the latter from $U_{m-1}(\cos\theta) = \sin(m\theta)/\sin\theta$)
collapse each Egypt term to a single expression at the seed~$z$:
\begin{equation}\label{eq:termcollapse}
\text{Term}_k
= \frac{U_{m-1}(z)}{
  T_{\lceil k/2\rceil m}(z)\;\bigl(
    U_{(\lfloor k/2\rfloor+1)m-1}(z)
  - U_{\lfloor k/2\rfloor m - 1}(z)\bigr)}.
\end{equation}

Since $z$ is a small rational number (while $x = T_m(z)$ can be
astronomically large), this expresses every term as a ratio of
Chebyshev polynomials evaluated at a manageable argument.

\begin{example}
For $n = a^2 + 4$ ($a$ odd): $m = 3$, $z = (a^2+2)/2$,
$U_2(z) = (a^2+1)(a^2+3)$.
\begin{align*}
\mathrm{Term}_1 &= \frac{1}{T_3(z)}
  = \frac{2}{(a^2+2)(a^4+4a^2+1)}, \\[3pt]
E_0 &= \frac{a(a^2+3)}{a^2+1}, \qquad
n - E_0^2 = \frac{4}{(a^2+1)^2}.
\end{align*}
Both the approximation and its error are explicit polynomials in~$a$.
\end{example}

%% ====================================================================
\section{Regulator interpretation}
%% ====================================================================

The regulator $R = \log\varepsilon$ measures the ``size'' of the
fundamental unit in $\mathbb{Q}(\sqrt{n})$.
Theorem~\ref{thm:tanh} gives it a direct approximation-theoretic meaning:

\begin{itemize}
\item $E_k \to \sqrt{n}$ as $\tanh((k{+}1)R/2) \to 1$,
  i.e., as $(k{+}1)R/2 \to \infty$.
\item The interval width $\sqrt{n}/\sinh((k{+}1)R)$
  decays as $2\sqrt{n}\,\varepsilon^{-(k+1)}$.
\item For small regulators (easy Pell equations), convergence is slow.
  For large regulators (hard Pell equations), even $E_0$ is an
  excellent approximation: the relative error is $O(\varepsilon^{-1})$.
\end{itemize}

In particular, the Egypt series is most useful precisely when the
Pell equation is hardest to solve---the regulator that makes the
equation difficult also guarantees rapid convergence once the
solution is found.

\section*{Acknowledgments}

Computations were performed using Wolfram Mathematica.
The Egypt series originated in the author's
\texttt{egypt} package~\cite{egypt-repo}.

\bibliographystyle{plain}
\bibliography{references}

\end{document}
